%& -shell-escape
% !TeX root = thesis-abstract.tex
% !TeX encoding = UTF-8
% !TeX program = XeLaTeX

\documentclass{.local/lib/classes/ndvr-super-article-standalone}

\begin{document}
    \begin{titlepage}
        \begin{annotation-item}%
            {Метод поиска нечетких дубликатов видео}
            {Никитин И. К., научный руководитель Лукин В. Н.}
            Доклад посвящён поиску похожих видео. 
            Введено понятие «нечетких дубликатов» видео. 
            Нечетким дубликатом называют объект, 
            частично совпадающий 
            с~другим объектом подобного класса. 
            Рассмотрена актуальность проблемы 
            поиска нечетких дубликатов видео. 
            Построена обобщённая модель поиска по видео. 
            Предложен метод поиска и сравнения нечетких дубликатов, 
            который с помощью сегментации видео по временной оси, 
            сводит задачу к поиску и сравнению временных рядов. 
            Представлен декларативный язык потоковой фильтрации видео, 
            с~помощью которого реализована сегментации. 
            Сформулирован алгоритм индексации 
            и поиска нечётких дубликатов видео. 
            Проведено сравнение с аналогами.
        \end{annotation-item}

        \begin{annotation-item}%
            {Near-duplicate Video Retrieval Method}
            {Nikitin I. K., scientific adviser Lukin V. N.}
            \begin{english}
                The presentation focuses 
                on finding of similar videos. 
                The notion of «near-duplicates» is introduced. 
                A near-duplicate is an object that approximates 
                another object of the same class. 
                The presentation contains 
                near duplicate problem statement 
                and a generalized model of the video search. 
                Also a method of searching and comparing 
                near duplicate videos is proposed.  
                The method reduces near duplicate videos retrieval 
                to time series indexation. 
                This is achieved through a temporal video segmentation. 
                Also a streaming video filtering 
                declarative language is presented. 
                It helps to implement a video segmentation. 
                A final algorithm for near duplicate videos 
                indexation and retrieval is formulated 
                and compared with analogues.
            \end{english}

        \end{annotation-item}
    \end{titlepage}
    \newpage
\end{document}
